%! Author = Sujal Singh
%! Date = 4/9/24

% Preamble
\documentclass[11pt,t]{beamer}
\title[Moulding of Thermosets]{\textbf{\Large Moulding of Thermosets}\vspace*{-10pt}}
\author[Sujal, Divyanshi, Priyanshu \ldots]{%
    \textbf{Group 9:}\\%
    \(|\) Sujal Singh \(|\) Divyanshi Panchal \(|\) Priyanshu Raj \(|\)\\\(|\) Chitransh Koshta \(|\) Prashant
    Pulkit \(|\)\\[10pt]%
    \textbf{Submitted to:}\\Mr. Rajendra Prasad\vspace*{-20pt}
}
\date[Chitransh, Prashant]{\textbf{Enrollment Numbers:}\\0\{41-45\}19051723}

\usetheme{Boadilla}
\setbeamertemplate{navigation symbols}{}
\setbeamertemplate{frametitle continuation}[from second][]
\setlength{\leftmarginii}{12pt}
\setlength{\leftmargini}{12pt}

% Packages
\usepackage{amsmath}
\usepackage[edges]{forest}
\usepackage{graphicx}
\usepackage{textcomp}
\usepackage{chemfig}

\setbeamertemplate{caption}[numbered]

% Document
\begin{document}
    \begin{frame}{Manufacturing Presentation (Group 9)}
        \\[5pt]
        \begin{center}
            %! suppress = FileNotFound
            \includegraphics[width=70pt]{logo}
        \end{center}\vspace*{-15pt}
        \maketitle
    \end{frame}

    \begin{frame}[allowframebreaks]{Index}
        \begin{enumerate}
            \large
            \item Thermosets
            \item Examples \& Applications of Thermosets
            \item Moulding Process
            \item Compression Moulding
            \item Transfer Moulding
        \end{enumerate}
    \end{frame}

    \begin{frame}[allowframebreaks]{Thermosets}
        \begin{enumerate}
            \item \textbf{Plastics} which are set under the application of \textbf{heat} and/or \textbf{pressure}.
            \item Thermosets harden when they are heated, if heated further they will \underline{breakdown chemically
            and lose their properties}.
            \item The process is \textbf{not reversible}, hence thermosets \underline{can not be recycled}.
            \item They consist of \textbf{3D network structures} based on strong covalent bonds to form rigid solids
            \& linear molecular chains are connected by weak secondary bonds.
        \end{enumerate}

        \framebreak

        \begin{minipage}[t]{\textwidth}
            \begin{center}
                \begin{figure}
                    \scalebox{0.6}{
                        \begin{forest}
                            forked edges,
                            for tree={draw,font=\strut\footnotesize\sffamily\em}
                            [Polymeric materials,s sep=150pt
                            [Elastomers]
                            [Plastics,s sep=50pt
                            [Thermoplastics
                            [Comodity Plastics
                            [Polyethylene,draw=none[Polypropylene,draw=none[Polystyrene,draw=none[PVC,draw=none]]]]]
                            [Engineering Plastics
                            [Polyamides,draw=none[Acetals,draw=none[Cellulosics,draw=none[Polycarbonate,draw=none,[
                            Others
                            ,draw=none]]]]]]]
                            [Thermosetting
                            [Commodity Plastics[Phenolics,draw=none[Unsaturated polyesters,draw=none[Ureas,draw=none]]]]
                            [Engineering Plastics[Silicones,draw=none[Urethanes,draw=none[Melamines,draw=none[Others,
                            draw=none]]]]]]]]]
                        \end{forest}
                    }
                    \caption{Classification of the different polymeric materials}\label{fig:fig-1}
                \end{figure}
            \end{center}~\\[10pt]
        \end{minipage}


        \framebreak

        \begin{itemize}
            \item \textbf{Epoxies}
            \begin{itemize}
                \item \textbf{Characteristics:} Excellent combination of mechanical properties and corrosion
                resistance, dimensionally stable, good adhesion, relatively inexpensive, good electrical properties.
                \item \textbf{Applications:} Electrical mouldings, sinks, adhesives, protective coatings, used with
                fiberglass laminates.
            \end{itemize}
            \item \textbf{Phenolic}
            \begin{itemize}
                \item \textbf{Characteristics:} Excellent thermal stability to over 150\textdegree C, may be
                compounded with a large number of resins, fillers, etc., inexpensive.
                \item \textbf{Applications:} Motor Housing, telephones, auto distributors, electrical fixtures.
            \end{itemize}
            \item \textbf{Polyester (PET or PETE)}
            \begin{itemize}
                \item \textbf{Characteristics:} One of the toughest of plastic films, excellent fatigue and tear
                strength, and resistance to humidity acids, greases, oils and solvents.
                \item \textbf{Applications:} Magnetic recording tapes, clothing, automotive tire cords, beverage
                containers.
            \end{itemize}
        \end{itemize}
    \end{frame}

    \begin{frame}[allowframebreaks]{Moulding Process}
        \begin{itemize}
            \item This process is used mostly for thermo-setting polymers.
            \item In compression moulding, the raw material is placed inside the mould in semi-solid (i.e., as
            granules, or a single piece called a plug).\ The mould is heated and closed using pressure, and the plastic
            flows to fill the cavity.
            \item If the part shape is more complex, transfer moulding may be used.\ Here, the charge are placed in a
            heated cylinder till they are soft, a hole at the bottom of the cylinder is connected to the die cavity
            by a sprue.
            \item A plunger pushes the semi-solid plastic into the die through the sprue, using high pressure.
        \end{itemize}
    \end{frame}

    \begin{frame}[allowframebreaks]{Compression Moulding}
        \begin{figure}
            \vspace*{-10pt}
            \begin{center}
                \includegraphics[width=200pt]{compression_molding}
            \end{center}
            \caption{Compression Moulding}\label{fig:fig-2}
            \vspace*{-10pt}
        \end{figure}
        \begin{itemize}
            \item A compression mould is made of two halves with one each being connected to the platens of the press.
            \item The mould is electrically heated to maintain the required temperature.
            \item Material is placed in the mould, and it is closed with a hydraulic cylinder, or toggle clamp.
        \end{itemize}

        \framebreak

        ~\\
        \begin{minipage}{0.6\textwidth}
            \begin{itemize}
                \item The pressure maintained on the material is of the order of 14 to 40 MPa of moulding area.
                \item As the material comes in contact with the heated mould surface, it softens and fills the entire
                cavity and at the same time intimates the chemical reaction which cures the part.
                \item Cure time is determined by the thickest cross-section, mould temperature, material type and grade.
                \item The most widely used plastic is phenol-formaldehyde, commonly known as \textbf{`Bakelite'}.
            \end{itemize}
        \end{minipage}
        \begin{minipage}{0.385\textwidth}
            \begin{center}
                \begin{figure}
                    \chemfig[atom style={scale=0.45}]{%
                        *6(-(-CH_2-)=-(-CH_2-[::-60]*6(=-(-CH_2-*6(-(-CH_2-[::60]*6(=-(-CH_2)=-(-OH)=-))=-(-OH)=-(
                        -CH_2)=
                        ))=-
                        (-CH_2)=(-OH)-))=(-OH)-(-CH_2)=)
                    }
                    \caption{\smallsize Bakelite}\label{fig:fig-3}
                \end{figure}
            \end{center}
        \end{minipage}
    \end{frame}

    \begin{frame}[allowframebreaks]{Transfer Moulding}
        \begin{figure}
            \vspace*{-10pt}
            \begin{center}
                \includegraphics[width=160pt]{transfer_molding}
            \end{center}
            \caption{Transfer Moulding}\label{fig:fig-4}
            \vspace*{-10pt}
        \end{figure}
        \begin{itemize}
            \item Transfer moulding is a manufacturing process where casting material is forced into a mould.
            \item This is similar to compression moulding but the main difference is that in compression the mould is
            enclosed whereas in transfer the mould is open to the fill plunger resulting in higher dimensional
            tolerances.
            \item Thermoset polymer and epoxy resins are some materials which are used in transfer moulding.
        \end{itemize}

        \framebreak

        \textbf{Process:}\\
        \begin{itemize}
            \item The pre-heated, uncured compound is placed in the transfer pot.
            \item The mould is closed up and under hydraulic pressure, the moulding compound is forced through a
            small hole into the cavity.
            \item The mould is held closed while the moulding compound is cured (thermosets) or cooled (thermoplastics).
            \item The mould is then split to free the product, with the help of ejector pins.
            \item Any flash or sprue material is trimmed off.
        \end{itemize}

    \end{frame}

    \begin{frame}{References}
        \begin{itemize}
            \item Manufacturing Technology: Foundry, Forming and Welding -- P N Rao
            \item \href{https://wikipeida.com}{\color{blue}Wikipedia}
        \end{itemize}
    \end{frame}

    \begin{frame}{The End}
        \begin{center}
            \vfill
            \textbf{\Huge Thank You.}
            \vfill
        \end{center}
    \end{frame}
\end{document}
